\documentclass[11pt]{article}
\usepackage[margin=1in]{geometry}
\usepackage{amsmath,amssymb,amsthm,bm}
\usepackage{booktabs,tabularx,longtable}
\usepackage{enumitem}
\usepackage{natbib}
\usepackage[hidelinks]{hyperref}
\setlength{\emergencystretch}{2em}
\newtheorem{definition}{Definition}
\newtheorem{theorem}{Theorem}
\newtheorem{proposition}{Proposition}
\newtheorem{corollary}{Corollary}

\title{A Scalar-15 Source Preimage and Its Renormalization Boundary\\
\large Tau Core Technical Paper VII-C}
\author{Jozsef Olcsak}
\date{Last revised: 18 September 2026}

\begin{document}
\maketitle

\begin{abstract}
We isolate the scalar-15/type-II branch previously recorded as a downstream
handoff in the joint-terminal paper. Representation type alone does not force
a nonzero field: with positive stiffness and no occupied dual source, the zero
section remains the unique local minimum. In a declared recovery-complete
one-carrier class, however, a symmetric two-Higgs seed edge is the minimal
preimage of a Majorana terminal plus a positive Higgs-dependent bosonic mode.
A faithful base realization and one positive matching row then produce the
cubic incidence, its source-neutral quartic completion and an induced occupied
scalar branch. Radiative closure fixes the invariant operator support but not
its coefficient ray. A covariant spectral-tensor operator selects one
conditional boundary ray, while exact counterfamilies show that the unchanged
source and a common unit do not select a unique RG trajectory. We separate
class-relative source ownership, the ratio $M_*/E_P=1/\sqrt{24\pi}$ with one
CODATA calibration, and the standard exact functional-RG evaluation law.
Physical type-II occupation, full component beta functions and unrestricted
parent realization remain open.
\end{abstract}

\paragraph{Inherited morphology terminology.}
The scalar branch is evaluated only after the upstream base carrier and its
base-readable seed/load pattern select a complete \emph{stabilized
morphological response configuration} $M_\tau^\star$ (short name:
morphological body).  This paper does not infer the body from the scalar
terminal and does not identify atemporal selection with literal formation or
observer time.
Under MRC-DEF2, these parent-side objects are typed pregeometric relational
candidates, not ordinary bodies already embedded in recovered spacetime;
ordinary matter objects occur only downstream of the supplied body.

\begin{center}
\fbox{\begin{minipage}{0.92\linewidth}\small
\textbf{Inputs.} Papers I--VI supply the body-first architecture, typed source
support, recovery boundary, conditional parent-realization language, temporal
descent and quantum terminal. Paper VII-A owns joint gravity--matter--Abelian
gauge assembly; Technical Paper VII-B owns the rank-seven P3 packet.

\textbf{Results owned here.} The scalar-15 occupation no-go, induced
occupation theorem, typed seed-edge and positive matching-row origin, source
rank of the scalar boundary packet, radiatively closed invariant support,
spectral-tensor boundary ray, class-relative ownership closure, common-scale
ratio and exact functional-RG evaluation boundary.
\end{minipage}}
\end{center}

\begin{center}
\fbox{\begin{minipage}{0.92\linewidth}\small
\textbf{Claim boundary.} The branch is selected only inside the declared
recovery-complete, source-generated and source-irredundant P4/type-II class.
Observed neutrino mass does not select type II over type I or III. The SI mass
uses measured constants, and the exact functional flow does not supply the
unevaluated component beta polynomials or prove physical branch occupation.

\textbf{Inherited formation basis.} The scalar branch is downstream of a
stable selected nonneutral body.  The base is the protected zero-source
response structure and the seed is a response-distinct base-relative source
class of one joint law.  Fixing a scalar terminal branch cannot supply the
physical zero-source anchor or remove the upstream base--seed
refactorization freedom.
\end{minipage}}
\end{center}

\tableofcontents
\clearpage

\section{Relation to established particle physics}

\paragraph{Inherited observer co-descent convention.}
The scalar/RG terminal is evaluated only after operational observer support
has formed: an occupied body-side carrier, regular local rank-four descent,
stable observer--source quantizer and nonzero record/effect co-produce the
observer and its accessible 4D world.  The RG construction neither creates
that observer nor defines an independent channel layer; it specializes an
already occupied terminal branch.

\paragraph{Imported descriptor-migration and response boundary.}
This specialization consumes the typed continuous descriptor and operational
record
\begin{equation}
 \Xi_{OS}^{\rm cont}=\Phi_{OS}[M_\tau],
 \qquad D_{OS}^{\rm op}=Q_{OS,\delta}(\Xi_{OS}^{\rm cont}),
 \label{eq:viic-cont-operational}
\end{equation}
without identifying them.  A historical scalar or RG descriptor $L$ can be
reused only if $L=U\circ\Xi_{OS}^{\rm cont}$, equivalently
$\ker D\Xi_{OS}^{\rm cont}\subseteq\ker DL$ on a regular linearized patch;
equal rank alone does not establish factorization.  The upstream post-body
response is fixed conditionally by $\Gamma=-H^{-1}B$, with the SBTCL106
specialization $\Gamma=J_{OS}$ and $C-B^*H^{-1}B=0$.  An identical-reduct
$\lambda$ counterfamily proves that positivity, common mediation and Schur
neutrality do not select the vertical source jet.  These are dependency and
claim-boundary statements only: this paper does not derive $Q_{OS,\delta}$,
terminal calibration, a galactic $q_R(R)$, the relevant physical $J_{OS}$
component or its Nature occupation.

Scalar-triplet/type-II mechanisms for Majorana neutrino mass and the relevant
unified-group representation theory are established particle-physics
constructions \cite{MaggWetterich1980,SchechterValle1980,Slansky1981}. General
renormalizable field theories already possess systematic beta-function and
exact functional-flow formalisms \cite{MachacekVaughn1985,Wetterich1993}.
This paper does not claim those mechanisms or equations as Tau-specific. Its
question is whether a single typed seed--base source law can reconstruct the
branch's representation, occupation, boundary support and common calibration
without independently fitted terminal gains.

\section{Type does not imply occupation}

Let $V\simeq\mathbb C^5$ be the Higgs carrier and
$\Delta\in\operatorname{Sym}^2V$ the scalar-15 coordinate. After the body and
an occupied Higgs coordinate $H$ are frozen, the first local analytic action
has the form
\begin{equation}
\begin{split}
\Gamma_{15}[\Delta;H]={}&m_{\rm eff}^2(H)\lVert\Delta\rVert^2
+\lambda_1\lVert\Delta\rVert^4
+\lambda_2\operatorname{Tr}(\Delta^\dagger\Delta\Delta^\dagger\Delta)\\
&-2\operatorname{Re}\langle J_{15}(H,\psi),\Delta\rangle+\cdots,
\end{split}
\label{eq:localaction}
\end{equation}
where $m_{\rm eff}^2=M_\Delta^2+\kappa_{H\Delta}H^\dagger H$. An unbroken
$SU(5)$-covariant source obeys
$\operatorname{Hom}_{SU(5)}(15_H,\mathbf1)=0$, so it cannot provide a
constant nonsinglet load. A portal changes the Hessian but its first variation
vanishes at $\Delta=0$. If $m_{\rm eff}^2>0$, the quartic form is nonnegative
and no dual source or zero-excluding topology is present, the zero section is
the unique local minimum. Hence
\begin{equation}
\boxed{B_\tau^{\rm phys}+s_U+\text{current P3 law}
\not\Longrightarrow\Delta_*\ne0.}
\label{eq:occupationnogo}
\end{equation}

\begin{proposition}[Induced scalar-15 occupation]
If $H\ne0$, $K_{15}>0$ and the source owns a nonzero covariant incidence
$\mu_\Delta\langle(H\odot H)^\flat,\Delta\rangle+\mathrm{h.c.}$, then the
unique small-source branch satisfies
\begin{equation}
\boxed{\Delta_*=K_{15}^{-1}\bar\mu_\Delta(H\odot H)^\flat
+O(\lVert\mu_\Delta H\odot H\rVert^3).}
\label{eq:inducedoccupation}
\end{equation}
\end{proposition}
On the electroweak triplet component this gives
$v_\Delta\simeq\mu_\Delta v_H^2/M_{\Delta,\rm eff}^2$ and
$m_\nu\simeq Y_\Delta v_\Delta$. This proves neither $\mu_\Delta\ne0$ nor the
Yukawa map $Y_\Delta$.

\section{Typed seed edge and positive base realization}

The seed does not contain two Higgs particles. It contains an unordered
arity-two typed incidence
\begin{equation}
e_{HH\Delta}=(V_H,V_H;S_{15})/S_2,
\qquad H\mapsto H\odot H.
\label{eq:typededge}
\end{equation}
Let the enriched base realize it by an equivariant map
$C_{15}:\operatorname{Sym}^2V\to15_H$. Since both sides are the same
irreducible representation, Schur's lemma gives
$C_{15}=\mu_\Delta I_{15}$. For $K_{15}>0$, define
\begin{equation}
R_{15}=K_{15}^{1/2}\Delta-K_{15}^{-1/2}C_{15}(H\odot H).
\label{eq:matchingrow}
\end{equation}
Its once-counted positive action is
\begin{equation}
\begin{split}
\frac12\lVert R_{15}\rVert^2={}&
\frac12\langle\Delta,K_{15}\Delta\rangle
-\operatorname{Re}\langle C_{15}(H\odot H),\Delta\rangle\\
&+\frac12\langle C_{15}(H\odot H),
K_{15}^{-1}C_{15}(H\odot H)\rangle.
\end{split}
\label{eq:matchingaction}
\end{equation}
The middle term is the cubic incidence and the final term is its required
source-neutral completion, not a separately fitted Higgs interaction.
Variation gives $K_{15}\Delta_*=C_{15}(H\odot H)$.

Inside the frozen class requiring a renormalizable nonsinglet Majorana
mediator, one positive Higgs-dependent bosonic repair mode and one-carrier
source irredundancy, $15_H$ and Eq.~\eqref{eq:typededge} are the minimal common
carrier. Type-I and type-III fermionic mediators fail the bosonic one-carrier
requirement but remain alternatives outside this class. Neutrino mass alone
therefore does not select the branch.

\section{The cubic edge is not the complete RG packet}

The matching action contains stiffness, cubic incidence and a neutral $H^4$
completion, but no $H^2\Delta^2$ portal and no breaking descriptor that splits
the unified components. The weakest field-dependent stiffness is
\begin{equation}
K_{15}(H,\Sigma)=K_0+\kappa_{H\Delta}(H^\dagger H)I_{15}
+a_\Sigma\rho_{15}(\Sigma)+\cdots.
\end{equation}
Portal and breaking coefficients can vary while the cubic edge remains fixed;
this is an exact source-rank obstruction.

Assume local $SU(5)$ covariance, power-counting renormalizability, radiative
closure and no additional forbidding symmetry. A class-minimal support is
\begin{equation}
\begin{aligned}
I_{p1}&=(H^\dagger H)\operatorname{Tr}(\Delta\Delta^\dagger),&
I_{p2}&=H^\dagger\Delta\Delta^\dagger H,\\
I_b&=\langle\Delta,\rho_{15}(\Sigma)\Delta\rangle,&
I_{41}&=[\operatorname{Tr}(\Delta\Delta^\dagger)]^2,\\
&&I_{42}&=\operatorname{Tr}[(\Delta\Delta^\dagger)^2].
\end{aligned}
\label{eq:invariantsupport}
\end{equation}
This fixes operator directions, not the boundary vector
$c_I=(\kappa_1,\kappa_2,a_\Sigma,\lambda_{41},\lambda_{42})$. A common unit
cannot remove the remaining projective ratios. The finite selector
\begin{equation}
\mathcal L_I(c)=\tfrac12\langle c,G_Ic\rangle-\langle j_I,c\rangle,
\qquad c_I^*=G_I^{-1}j_I,
\end{equation}
is complete only after $G_I$ and $j_I$ are source-derived.

\section{Spectral-tensor boundary ray}

On $\mathcal H_{15}=\mathbb C\oplus V\oplus\bar V$, define
\begin{equation}
\mathcal D_{15}(H,\Delta,\Sigma)=
\begin{pmatrix}
0&H^\dagger&H^T\\ H&\Sigma&\Delta\\
H^*&\Delta^\dagger&\Sigma^T
\end{pmatrix}.
\label{eq:spectraloperator}
\end{equation}
It transforms by unitary conjugation under $SU(5)$. Its cubic moment contains
the breaking and charged incidences in one word,
\begin{equation}
\operatorname{Tr}\mathcal D_{15}^3
=6I_{H\Sigma}+6I_{HH\Delta}+3I_b+2I_{\Sigma3}.
\end{equation}
The canonical quartic completion
$\mathcal Q_4=\operatorname{Tr}\mathcal D_{15}^4
+(\operatorname{Tr}\mathcal D_{15}^2)^2$ gives at $\Sigma=0$
\begin{equation}
\boxed{\mathcal Q_4=24I_{H4}+16I_{p1}+8I_{p2}+4I_{41}+2I_{42}.}
\label{eq:spectralray}
\end{equation}
All directions in Eq.~\eqref{eq:invariantsupport} are present, and the portal
and scalar-quartic rays are each $2{:}1$. Equal weighting of single-trace and
tensor-square terms remains a completion premise, not a narrow-parent theorem.

\section{Common scale and metrological boundary}

The inherited common-length completion gives
\begin{equation}
\ell_H=\sqrt{48\pi}\,\ell_P,
\qquad
\boxed{M_*=\frac{\sqrt2\hbar c}{\ell_H}
=\frac{E_P}{\sqrt{24\pi}}
=\frac{\overline M_Pc^2}{\sqrt3}.}
\label{eq:commonscale}
\end{equation}
Using 2022 CODATA values \cite{CODATA2022} yields
$M_*=1.406035(16)\times10^{18}\,\mathrm{GeV}$. The Tau-specific statement is
the dimensionless ratio and absence of an independent scalar-15 length. The SI
number is a one-anchor calibration, not a prediction of $G$, $\hbar$ or $c$.

\section{Exact functional flow and its evaluation boundary}

Let $P$ be the positive background-covariant Laplace-type operator including
declared gauge-fixing and ghost blocks. The smooth spectral regulator
$R_k(P)=P/[\exp(P/k^2)-1]$ enters the exact effective-average-action flow
\cite{Wetterich1993}
\begin{equation}
\boxed{\partial_t\Gamma_k=\frac12\operatorname{STr}
[(\Gamma_k^{(2)}+R_k)^{-1}\partial_tR_k],\qquad t=\log k.}
\label{eq:wetterich}
\end{equation}
Gauge sectors must also satisfy the regulator-modified Slavnov--Taylor/BV
identity, whose insertion vanishes as $k\to0$. For a radiatively and
BRST-complete independent basis $\{\mathcal O_I\}$,
\begin{equation}
\boxed{\beta^I=(G^{-1})^{IJ}
\langle\mathcal O_J,\partial_t\Gamma_k\rangle,
\quad G_{IJ}=\langle\mathcal O_I,\mathcal O_J\rangle.}
\label{eq:betaprojection}
\end{equation}
These equations define an exact functional evolution and coordinate
projection after field content, scheme and matching are declared. They are
standard QFT machinery, not a Tau derivation of the Wetterich equation, and do
not evaluate the full component beta polynomials. Residual regulator
dependence in a truncation is an uncertainty diagnostic.

For one active complex $15_H$, representation theory fixes
$\Delta b^{15_H}=(7/6,7/6,7/6)$. This changes common running but not relative
one-loop gauge slopes. A unique trajectory still requires portals, scalar
quartics, Yukawa/cubic data, component thresholds and breaking-scale matching.
Equal-upstream choices of undetermined data can produce different
stationarity verdicts, so a chosen trajectory is not a Tau prediction.

\section{Four data classes that must not be conflated}
\label{sec:fourclasses}

A representation label and one cubic incidence do not define a physical RG
trajectory.  The finite obstruction can be stated without a parameter scan.
For a radial spectral potential
\begin{equation}
 V(x)=a_2x^2+a_3x^3+a_4x^4,
 \label{eq:radialpotential}
\end{equation}
four logically different data are required.

\begin{enumerate}[label=(\roman*)]
\item A \emph{moment character} fixes ratios such as
  $a_3^2/(a_2a_4)$ and the relative weight of single-trace and tensor-square
  invariants.  A common overall action unit cannot fix these ratios.
\item \emph{Pointed occupation} excludes the zero section.  Every unpointed
  polynomial of the form \eqref{eq:radialpotential} has $V'(0)=0$; covariance
  and positivity alone therefore do not force a nonzero occupied branch.
\item A \emph{physical scale} maps one dimensionless parent gap to a measured
  mass unit.  Rescaling all parent gaps preserves the dimensionless packet.
\item A \emph{renormalization map} specifies fields, gauge fixing, ghosts,
  scheme, thresholds and matching.  Finite coupling redefinitions change the
  coordinate representation of a beta vector even when observables agree.
\end{enumerate}

This separation prevents four common overclaims: normalization is not a
moment selector; spontaneous occupation is not implied by representation
content; a dimensionless scale ratio is not an SI prediction; and an exact
functional equation is not an evaluated component flow.

\begin{proposition}[Finite four-class obstruction]
The current scalar-15 edge, common unit and invariant support do not by
themselves determine a unique physical RG initial-value problem.
\end{proposition}
\begin{proof}
Multiplication by a common unit leaves $a_3^2/(a_2a_4)$ invariant.  Equation
\eqref{eq:radialpotential} has zero first derivative at the origin.  A common
rescaling of all dimensionful gaps leaves every normalized source relation
unchanged.  Finally, under $g'=g+cg^2$ a coordinate beta function transforms
as $\beta_{g'}=(1+2cg)\beta_g$.  These four independent transformations leave
the previously reconstructed packet intact while changing respectively its
moment shape, occupation verdict, SI scale or flow coordinates.
\end{proof}

\section{Conditional common origin of the four classes}
\label{sec:commonorigin}

The four classes need not be four unrelated additions.  One conditional
source architecture can own all four, provided its transport to the
scalar-15 operator is proved.

First, an oriented primitive response character supplies the moment sequence
\begin{equation}
 e^x=1+x+\frac{x^2}{2}+\frac{x^3}{6}+\frac{x^4}{24}+\cdots .
 \label{eq:momentcharacter}
\end{equation}
Transport of this one character to $\mathcal D_{15}$ fixes relative spectral
moments; its two-copy multiplicative character supplies the tensor-square
sector.  Reversal averaging cannot be imposed before the odd cubic has
selected an oriented branch, because doing so would delete the very incidence
that induces the scalar response.

Second, a pointed seed orientation may activate an already selected Higgs
body coordinate,
\begin{equation}
 V_H(h)=-k_Hh^2+\lambda_Hh^4,
 \qquad h_*^2=\frac{k_H}{2\lambda_H}>0.
\end{equation}
The positive matching row of Eq.~\eqref{eq:matchingrow} then yields
\begin{equation}
 \boxed{\Delta_*=\frac{\mu_\Delta}{K_\Delta}h_*^2
 =\frac{k_H\mu_\Delta}{2K_\Delta\lambda_H}\ne0.}
 \label{eq:pointedoccupation}
\end{equation}
This is not an external gauge-charged linear source: the seed selects an
oriented orbit, the body occupies the Higgs coordinate, and the covariant
two-Higgs edge induces the scalar coordinate.

Third, write the physical operator as
\begin{equation}
 \mathcal D_{15}^{\rm phys}=M_*\widehat{\mathcal D}_{15}.
 \label{eq:physicalgap}
\end{equation}
This identifies where the common scale lives while retaining the metrological
boundary stated in Section~\ref{sec:scalar15handoff}.

Fourth, for observer resolution $k$ the covariant spectral family
\begin{equation}
 P_k(\mathcal D_{15})=\mathbf 1_{[-k,k]}(\mathcal D_{15}),
 \qquad
 \Gamma_k=\operatorname{Desc}_{P_k}\Gamma
 \label{eq:spectraldescent}
\end{equation}
is nested and gauge covariant:
$P_{k_1}P_{k_2}=P_{k_1}$ for $k_1\le k_2$, and
$P_k(W\mathcal DW^\dagger)=WP_k(\mathcal D)W^\dagger$.
Eigenvalue crossings therefore define source-owned threshold events rather
than thresholds inserted after a fit.  The smooth regulator used later is a
computational realization of the same resolution logic, not a new physical
coupling.

\begin{theorem}[Conditional four-origin assembly]
If the physical base--seed completion owns the oriented character transport,
the pointed Higgs cycle, one positive common gap and the covariant spectral
descent, then moment character, nonzero scalar occupation, relative scale and
renormalization prescription descend from one source architecture.
\end{theorem}
\begin{proof}
Equations \eqref{eq:momentcharacter}, \eqref{eq:pointedoccupation},
\eqref{eq:physicalgap} and \eqref{eq:spectraldescent} supply the four data
classes of Section~\ref{sec:fourclasses}.  Each construction is covariant on
its declared carrier, and none introduces a terminal-specific gain.  The
conclusion is conditional because occupation of this complete source packet
by the unrestricted physical parent is not derived.
\end{proof}

\section{Bosonic capacity is a viability test, not selection}

For active field-dependent modes $m_j^2=k_jr^2$, the one-loop logarithmic
weight is
\begin{equation}
 64\pi^2 B=\sum_j(-1)^{F_j}n_jk_j^2.
\end{equation}
At the frozen source point used by the audit, $g=1$ and
$y_u=0.9615522319$, the gauge--fermion deficit is
\begin{equation}
 D_{GF}=3y_u^4-\frac{171}{200}=1.70955948.
\end{equation}
A complex scalar $15_H$ has thirty real modes.  If all active modes share the
portal jet $\kappa_{H\Delta}$, their positive contribution is
\begin{equation}
 \Delta(64\pi^2B)_{15}=30\kappa_{H\Delta}^2,
\end{equation}
so one copy can reverse the sign whenever
\begin{equation}
 \boxed{|\kappa_{H\Delta}|>\sqrt{D_{GF}/30}=0.238716.}
 \label{eq:capacitybound}
\end{equation}
Equation~\eqref{eq:capacitybound} establishes raw source-scale capacity.  It
does not select the portal, prove a Coleman--Weinberg vacuum or remove
decoupling and threshold effects.

\section{Equal-upstream RG counterfamily}

The representation fixes the contribution of one active complex $15_H$ to
the GUT-normalized one-loop gauge coefficients,
\begin{equation}
 \boxed{\Delta b^{15_H}=\left(\frac76,\frac76,\frac76\right).}
 \label{eq:gaugeshift}
\end{equation}
Thus a degenerate complete multiplet changes common running but not relative
one-loop gauge slopes.  A physical trajectory additionally requires portals,
independent scalar quartics, Yukawa and cubic data, component thresholds and
breaking-scale matching.

A controlled equal-upstream audit makes the missing information visible.  At
one common active interval, two admissible portal choices give
\begin{center}
\begin{tabular}{@{}ccc@{}}
\toprule
$\kappa_{H\Delta}$ & $\min(\lambda+\beta_\lambda/4)$ & zero present?\\
\midrule
$0.80$ & $+0.071386$ & no\\
$1.20$ & $-0.660165$ & yes\\
\bottomrule
\end{tabular}
\end{center}
while all already frozen upstream data are held fixed.  This is a finite
non-identifiability witness: choosing either trajectory without an additional
source law would be a fit, not a Tau prediction.

\section{Exact finite stopping point}

The scalar-15 branch and its upstream renormalization interface close at five distinct levels once the 14 November 2026 three-pillar theorems (\texttt{U1}, \texttt{U2}, \texttt{F1}, \texttt{F3}, \texttt{HOMOG-K4}) are incorporated:
\begin{center}\small
\begin{tabularx}{0.96\linewidth}{@{}p{0.28\linewidth}p{0.34\linewidth}X@{}}
\toprule
Question & Proved result & Remaining boundary\\
\midrule
Source preimage & Unique in the declared recovery-complete type-II class &
Unrestricted physical scalar-15 branch occupation\\
Spectral/Clifford seed scale (\texttt{U1}, \texttt{F1}) & $\widehat{\Delta}_{K_4}=\frac{4}{3}P_\perp=G_{\mathrm{Cl}}$ unifies \texttt{C2} with \texttt{C1} at $r_{\mathrm{seed}}=1/4$; $\Pi^2=\Pi\implies\dot{r}\equiv 0$ &
Occupied non-triviality postulate (\texttt{POST-OCC-NU1})\\
Intensive UV ratios (\texttt{U2}) & $a_{\mathrm{UV}}^{-2}$ renormalizes $1/(4G_{\mathrm{ren}})$; $M_*/E_P=1/\sqrt{24\pi}$ and $\eta_{\mathrm{hor}}\equiv 1$ have \texttt{0.00e+00} UV drift &
One CODATA SI scale anchor\\
Simplicial \& 60-decade RG protection (\texttt{HOMOG-K4}, \texttt{F3}) & Local $S_4\times S_4\implies\kappa_{\mathrm{dir}}^{\mathrm{bare}}\equiv 0$; $\operatorname{Tr}(T_{\mathrm{def}})=0$ + Frullani tail universality &
Component scalar-15 beta functions and truncation error\\
Phenomenological trajectory & Exact covariant flow; counterfamily shows non-uniqueness from unit alone &
Portal, quartic, threshold and matching data\\
\bottomrule
\end{tabularx}
\end{center}

Crucially, the apparent upstream RG fragilities flagged in earlier audits (\texttt{F41, F44, F47}) are resolved without parameter tuning:
\begin{enumerate}
\item \emph{Intensive UV-invariance vs.\ extensive $1/(4G_{\mathrm{ren}})$ absorption (Theorem~\texttt{U2}):} While Srednicki entanglement entropy scales extensively as $c_S A_{\mathrm{AH}}/a_{\mathrm{UV}}^2$ under radial lattice refinement, the induced gravitational counterterm shifts $1/(4G_{\mathrm{ren}})=1/(4G_{\mathrm{bare}})+c_S/(\pi a_{\mathrm{UV}}^2)$ with the identical coefficient, leaving all intensive dimensionless ratios---including $M_*/E_P=1/\sqrt{24\pi}$, $\beta_{\mathrm{KMS}}|\kappa_{\mathrm{KH}}|/(2\pi c)=1$, and $\eta_{\mathrm{hor}}\equiv 1$---identically UV-invariant (\texttt{0.00e+00} drift).
\item \emph{Laplacian--Clifford identity, projector normalcy, and vanishing bare spectator stiffness (Theorems~\texttt{U1}, \texttt{F1}, \texttt{HOMOG-K4 H1--H2}):} Chung's normalized Laplacian on regular $K_4$, $\widehat{\Delta}_{K_4}=D^{-1/2}L_{K_4}D^{-1/2}=\frac{4}{3}P_\perp=G_{\mathrm{Cl}}$, eliminates the \texttt{C1/C2} norm bifurcation at $r_{\mathrm{seed}}=1/4$, while pure-state idempotence $\Pi^2=\Pi$ ($\dot{r}\equiv 0$) renders $\Gamma'_{1\text{-loop}}(1/4)<0$ normal to $\mathcal{P}_1(M_2(\mathbb{C}))$ alongside tangential $S^2$ Hessian $+4I_2\succ 0$. On periodic Freudenthal--Kuhn $K_4$ tilings, local $S_4\times S_4$ gauge invariance enforces $\kappa_{\mathrm{dir}}^{\mathrm{bare}}\equiv 0$ identically (\texttt{R4GATE-NATN1} closed), and defect trace-freeness $\operatorname{Tr}(T_{\mathrm{def}})=0$ protects the 60-decade flow (\texttt{F3}).
\end{enumerate}
This establishes a clean finite stopping point: neither normalizing the packet nor scanning RG portals can decide whether Nature occupies the scalar-15/type-II branch rather than an alternative neutrino/matter completion.

\section{Condensed three-level closure}
\label{sec:scalar15handoff}

The P3 spine fixes the common assembly language but does not itself choose a
non-Abelian particle completion.  A later P4 branch closes three previously
separate questions at three different logical levels.  We record that result
here because it is the first particle-sector specialization of the joint
source architecture; the underlying P4 reconstruction is not reclassified as
a theorem owned by Paper VII.

\subsection{Class-relative base--seed ownership}

Let \(\mathcal C_{15}^{\rm rec}\) denote the recovery-complete,
source-generated and source-irredundant P4/type-II class.  It contains the
recovered one-copy PIS--CAR packet, one occupied Higgs branch, the scalar-15
terminal and no independently sourced terminal-null operation ideal.  The
reverse source reconstruction identifies the pointed exterior \(V_5\) seed
grammar and a positive Hermitian/Hodge base realization.  The one-carrier
type-II selector then fixes the typed edge
\begin{equation}
 e_{HH\Delta}:\operatorname{Sym}^2V_5\longrightarrow S_{15},
 \qquad \dim\operatorname{Sym}^2V_5=15,
 \label{eq:scalar15edge}
\end{equation}
and faithful realization gives rank fifteen.  The unique shortest connected
alternating seed--Higgs word and its positive quartic completion imply
\begin{equation}
 \boxed{
 \mathcal C_{15}^{\rm rec}
 \Longleftrightarrow
 (\mathcal S_{\wedge,5},e_{HH\Delta};h_B,*_B)
 \Longrightarrow h_*\ne0,\quad\Delta_*\ne0 .}
 \label{eq:scalar15ownership}
\end{equation}
Thus charged character, occupation and quartic stabilization require no
independent source datum \emph{inside this recovery class}.  Equation
\eqref{eq:scalar15ownership} is a minimal-preimage reconstruction of the
selected branch.  It is not a theorem that the unrestricted physical parent
occupies type II rather than another neutrino/particle completion.

\subsection{Common spectral scale}

The same completion relates its Higgs length and scalar-15 mass to the common
Planck terminal:
\begin{equation}
 \ell_H=\sqrt{48\pi}\,\ell_P,
 \qquad
 M_*=\frac{\sqrt2\,\hbar c}{\ell_H}
 =\frac{E_P}{\sqrt{24\pi}}
 =\frac{\overline M_Pc^2}{\sqrt3},
 \label{eq:scalar15scale}
\end{equation}
where \(E_P=\sqrt{\hbar c^5/G}\) and
\(\overline M_Pc^2=E_P/\sqrt{8\pi}\).  With the 2022 CODATA constants
\citep{CODATA2022},
\begin{equation}
 \boxed{M_*=1.406035(16)\times10^{18}\ {\rm GeV}.}
 \label{eq:scalar15numeric}
\end{equation}
The Tau-specific statement is the dimensionless ratio
\(M_*/E_P=1/\sqrt{24\pi}\) and the absence of an independent scalar-15 length.
Equation~\eqref{eq:scalar15numeric} is a one-anchor metrological calibration,
not a first-principles prediction of the SI value of \(G\).

\subsection{Exact interacting functional law}

Lift the finite particle operator to the full graded four-dimensional carrier
and let \(P\) be the positive background-covariant Laplace-type operator,
including the declared gauge-fixing and ghost blocks.  The no-extra-profile
spectral scheme uses
\begin{equation}
 \boxed{R_k(P)=\frac{P}{\exp(P/k^2)-1}}
 \label{eq:scalar15regulator}
\end{equation}
in the exact effective-average-action flow \citep{Wetterich1993}
\begin{equation}
 \boxed{
 \partial_t\Gamma_k=\frac{1}{2}\operatorname{STr}
 \left[(\Gamma_k^{(2)}+R_k)^{-1}\partial_tR_k\right],
 \qquad t=\log k .}
 \label{eq:scalar15flow}
\end{equation}
For gauge sectors Eq.~\eqref{eq:scalar15flow} is solved together with the
regulator-modified Slavnov--Taylor/BV identity; the regulator insertion must
vanish as \(k\to0\).  If \(\{\mathcal O_I\}\) is a radiatively and
BRST-complete independent operator basis and
\(G_{IJ}=\langle\mathcal O_I,\mathcal O_J\rangle\) on the source-owned trace
domain, the component beta functions are the unique projection
\begin{equation}
 \boxed{
 \beta^I=(G^{-1})^{IJ}
 \left\langle\mathcal O_J,\partial_t\Gamma_k\right\rangle .}
 \label{eq:scalar15betaprojection}
\end{equation}
Equations~\eqref{eq:scalar15regulator}--
\eqref{eq:scalar15betaprojection} fix one gauge-covariant exact functional
law, threshold prescription and coordinate projection.  Computing the full
component polynomials remains an evaluation problem.  The regulator is an
analysis scheme, not a physical coupling, observer quantizer or additional
law requiring Nature-level selection; exact \(k=0\) observables are regulator
independent, whereas residual dependence of a finite truncation is an
uncertainty diagnostic.

The net closure is therefore deliberately three-tiered:
\begin{center}
\begin{tabular}{@{}ll@{}}
\toprule
Question & Status\\
\midrule
base--seed ownership & proved inside \(\mathcal C_{15}^{\rm rec}\)\\
absolute numerical scale & dimensionless ratio derived; SI value calibrated once\\
full beta law & exact functional flow and projection fixed\\
physical type-II occupation & empirical model-selection question\\
\bottomrule
\end{tabular}
\end{center}


\paragraph{Complete-base-modulus boundary.}
The scalar stiffnesses, matching data and threshold coordinates used in this
paper belong to the physical base-modulus component $\beta_B$ of the complete
source characteristic $\mathfrak c_{B,s}=(q_B,\beta_B)$. They cannot be
removed by a seed-only decoder; only $\ker D\mathfrak c_{B,s}$ is true
realization redundancy. RG counterfamilies that omit an independently
physical component of $\beta_B$ are therefore incomplete source models, not
proof that the corresponding modulus is gauge. Completeness of $\beta_B$ and
its occupation in Nature remain explicit open boundaries.

\section{Decision ledger}

\begin{center}\small
\begin{tabularx}{0.96\linewidth}{@{}p{0.42\linewidth}p{0.20\linewidth}X@{}}
\toprule
Statement & Status & Boundary\\
\midrule
Current P3 packet forces nonzero scalar-15 occupation & disproved & positive
zero-section counterbranch\\
Occupied Higgs plus a typed cubic induces $\Delta_*\ne0$ & proved locally &
requires the charged source edge\\
Declared one-carrier class selects the two-Higgs edge & proved conditionally &
type-I/III alternatives remain outside the class\\
One matching row owns stiffness, cubic and neutral quartic completion & proved
& does not own portal or component thresholds\\
Radiative closure fixes invariant operator support & proved conditionally &
coefficient values remain open\\
Common unit fixes a unique RG ray & disproved & projective ratios remain\\
Spectral tensor fixes a concrete boundary ray & constructed conditionally &
physical occupation is not derived\\
$M_*/E_P=1/\sqrt{24\pi}$ & class-relative derived ratio & SI number uses one
CODATA calibration\\
Exact functional RG and projection & standard exact evaluation law & component
calculation and truncation control remain work\\
Descriptor migration and post-body response & imported conditionally & fibre
factorization and physical source-jet occupation remain upstream\\
Unrestricted physical parent occupies type II & open & empirical or deeper
base--seed selection problem\\
\bottomrule
\end{tabularx}
\end{center}

\section{Scalar-branch occupation criterion}

The scalar-15 carrier is not occupied merely because its stiffness is
positive. Relative to the retained body its local component is
\[
 e_*=(C-B^\dagger A^{-1}B)^{-1}
      (k-B^\dagger A^{-1}j).
\]
The charged source edge and typed cubic used here provide the class-relative
nonzero loading required for the induced branch. Removing that edge gives the
positive zero-section control. This sharpens branch occupation inside the
declared type-II class but does not prove that the unrestricted physical
base--seed parent realizes that class or fixes its RG boundary data.

The scalar branch therefore inherits, but cannot prove, Paper IV's
operational physical-parent identity. Under the enriched one-source premise it
is evaluated on the seed-generated occupied quotient, and an inert ambient
summand does not modify the spectrum or flow. A non-null scalar/RG difference
at fixed complete base--seed key would instead identify a second source or an
incomplete upstream signature and falsify the one-source completion.

\subsection*{Global-law proof obligation}

The scalar-15 and RG branch is conditional on global body occupation.  The
source-fidelity/Hodge law suppresses every positive-Schur topology sector not
effectively loaded by the seed, but the present RG construction does not prove
that Nature owns that law or its component zero-jets.  A Nature-level spectrum
or running-coupling claim therefore requires independent source or empirical
evidence for the global parent law in addition to the branch-local occupation
and renormalization conditions derived here.

\subsection*{Unified-selection non-impact}

PD-USS1 conditionally selects a common state--transport carrier only after a
source-owned grading, reversible transport and actuality law are supplied.
PD-USS-N1 then proves that such selection does not force a nonzero
standard-orthogonal terminal effect. The updated $0/31$ P3 pair census is
therefore neither evidence for scalar-15 occupation nor a constraint on the
functional-RG boundary derived here. The scalar branch still requires its
own source load, occupied support and held-out spectrum/running-coupling test.

% BEGIN OBSERVER UPDATE 20260914
\section{Observer realization: inherited results and physical boundary}
\label{sec:observer-update-20260914}
The current observer programme distinguishes a conditional mathematical
construction from an occupied physical measuring system. The source order
remains base plus one universal seed, stabilized morphological response,
observer--source restriction, shared descriptor, and resolved terminal
readouts. Neither a detector nor an empirical residual constructs the parent
body backwards. Dynamics below concern downstream subsystems only.

The canonical observer-origin derivation, Sections 22--26, owns the following
results. The existing joint source action fixes static and frequency-dependent
response together. Variation of the BRAC potential couples to the field
square; a linear field coupling is a distinct comparator unless derived by
expansion about an occupied background. Fixed canonical means select a unique
coherent state in a finite positive harmonic packet only if there is no excess
fluctuation energy. This is a free/asymptotic preparation statement, not the
ground state of an interacting detector. For the existing coupled harmonic
pair, the joint ground is pure but either local subsystem is mixed, with its
covariance fixed by the complete joint Hamiltonian.

These are conditional derivations and internal numerical checks. The actual
observer--source context, physical preparation process, kinetic/energy
calibration and stable operational resolution remain unproved. The existing
ORA-T1 selector conditionally chooses stable quantizers when the occupied
context, admissible partition family and positive closure barrier are supplied;
it does not establish those physical premises. A computed ground state is not
a preparation mechanism, and a covariance or optimal abstract measurement is
not an implemented stable quantizer $Q_{OS}$.

\paragraph{Ownership and use.}
The central observer-origin note owns the new contact, state and covariance
derivations; the body-readout technical companion carries the consolidated
mathematics. Foundation Papers IV, V and VI retain physical-selection,
clock-descent and quantum-readout ownership, respectively. This paper imports
only the consequences within its scope. No observational dataset, baseline,
frozen coefficient, score or empirical conclusion is revised by this update.
The reproducibility supplement records the source hashes and the exact
conditional/physical distinction.


\paragraph{Consequence for this paper.}
These results do not select a particle spectrum, coupling constant, Higgs vacuum or gravitational terminal. Common source ancestry remains weaker than physical realization of those sectors.
% END OBSERVER UPDATE 20260914

% BEGIN STATE UPDATE 20260918
\section{Source identification and recovery}
The current observer audit separates a conditional forward construction from physical source selection. A frozen body, occupied observer incidence, preparation, complete instrument and stable resolution can determine records without selecting those inputs from Nature. Equal effects do not fix postmeasurement states, and symmetry within a chosen interaction family does not exclude other symmetric couplings. These results do not select a particle spectrum, coupling constant, Higgs vacuum or gravitational terminal. Common source ancestry remains weaker than physical realization of those sectors.

\paragraph{Dependency and claim ownership.}
Foundation VI and the body--readout technical companion own the finite instrument/access proof; Foundation IV and TD-0 retain physical source selection; Foundation V, the temporal synthesis and seed--body companion own the gradient/inertial distinction. This paper inherits only the stated specialization. Canonical assumptions, countermodels and calibration audits are in \url{https://github.com/tau-core-research/tau-core-theory/blob/main/docs/tau_rqm_instrument_to_record_bridge_001.md}. No physical seed, Nature occupation or Tau-exclusive observation is established by this update. Existing empirical endpoint definitions and scores are unchanged.
% END STATE UPDATE 20260918

% BEGIN PAIR READOUT UPDATE 20260927
\section{Conditional source preparation and time-resolved access}
The common-operator scalar--Dirac completion now supports an explicit
conditional chain from interaction vertices to pair transfer and joint
readout. Nonzero vertices do not select occupation. Relational coherence
can be generated from an occupied number-diagonal pair even while both
local mean fields vanish. In a declared resonant rotating-wave and Markov
loss model, finite-window emission retains this coherence information,
whereas an all-time count discards it. This is standard-compatible quantum
recovery, not a Tau-specific observational residual.
These results do not select a particle spectrum, coupling constant, Higgs vacuum or gravitational terminal. Common source ancestry remains weaker than physical realization of those sectors.
\paragraph{Dependency ownership and limits.}
The P4 common-source assembly owns the scalar/Yukawa benchmark and finite
preparation calculations; Foundation VI and the body--readout companion
inherit the joint-access mathematics. Foundation IV and TD-0 retain
physical source selection, and Foundation V retains clock descent.
The parent remains atemporal: the evolution above is a conditional
downstream 4D model, not parent memory or a derived morphogenetic time.
Physical initial occupation, a gauge-complete realization, detector
amplification, absolute calibration and stable observer-relative $Q_{OS}$
remain unselected. No empirical endpoint or score is changed.
The canonical derivation and reproducible numerical sources are linked at
\url{https://github.com/tau-core-research/tau-core-theory/blob/main/docs/tau_core_p4_nonabelian_particle_common_source_assembly_001.md}.
% END PAIR READOUT UPDATE 20260927

\section{Conclusion}

The scalar-15 branch now has a finite source ledger rather than a list of
freely chosen particle parameters. Representation type, activation
obstruction, induced occupation, missing source directions, radiatively
closed support and a concrete conditional boundary ray are separately
identified. The common scale is a dimensionless cross-terminal relation; its
SI value is calibrated rather than predicted. The functional RG supplies an
exact standard evolution equation once branch and scheme are declared, but it
does not evaluate the component beta functions.

The remaining scientific question is branch occupation: whether the physical
base--seed response realizes this scalar-15/type-II completion rather than
another neutrino or particle completion. That cannot be settled by normalizing
the packet, fitting a portal, or citing the exact RG equation.

\begingroup\small
\bibliographystyle{unsrtnat}
\bibliography{refs}
\endgroup
\end{document}
